% !TeX root = ./../main.tex

\section{Wyniki uzupełniające do rozdziału~\ref{chapter:classification}}\label{app:results:cls}
\agentinfo{
Utworzone w zadaniu PHD-TS-99 (2026-11-27) przy skracaniu rozdziału 8.
Tabele tab:cls:models, tab:cls:hw-results i tab:cls:pdu-recovery przeniesiono z podrozdziałów 8.3, 8.4.1 i 8.4.2 bez zmiany liczb; opis kalibracji stanami GHZ, parametry uczenia i procedurę dostrajania przeniesiono z podrozdziałów 8.3 i 8.4.1; tabela tab:cls:app:family jest nowa (liczby z plików porownania.csv obu konfiguracji eksperymentu porównawczego, rodzina glowna).
Rozdział 7 odwołuje się do wierszy M1, M2 i M6 tabeli tab:cls:models (poziom ibm_fez dla sześciu kubitów).
Objętość sekcji: około 1,8 strony, w tym cztery tabele około 1,3 strony (limit z zasad skracania: około 1,5); jeśli autor woli tabelę tab:cls:pdu-recovery w tekście głównym, rozdział 8 rośnie o około 0,3 strony.
Wiersz input tego pliku w chapters/appendix_results.tex dopisuje orkiestrator.
}

\paragraph{Modele i poziom ibm\_fez.}
Tabela~\ref{tab:cls:models} zestawia modele klasyfikatora uruchomione na urządzeniu (podrozdział~\ref{sec:cls:methodology}).
Poziom ibm\_fez wyznaczono z kalibracji stanami GHZ o $n$ od $2$ do $7$ kubitach, przygotowanymi na urządzeniu ibm\_fez i zmierzonymi z $8192$ powtórzeniami każdego obwodu, bez mitygacji błędów; wierność, wyznaczona metodą koherencji wielokwantowych~\cite{wei_verifying_2020}, wynosi od $0{,}985$ dla dwóch kubitów do $0{,}863$ dla siedmiu.
Trójkę $(\cohangle,\qphase,\qdep)$ dobrano tak, by symulacja obwodu przygotowującego stan, z szumem po początkowej bramce Hadamarda i po każdej bramce CNOT, odtwarzała zmierzoną wierność; jedna wierność nie wyznacza trzech parametrów jednoznacznie, a reguły podziału straty, przypisanej depolaryzacji i niewielkiemu odstrojeniu, dostępne źródła nie opisują.
% ŹRÓDŁO: torus_classifier/config.py (GRADED_NOISE_LEVELS, IBM_FEZ_GHZ6_CALIBRATION; komentarz: the inversion assigned the whole loss to depolarizing and a negligible coherent detuning) i torus_classifier/ibm_fez_ghz_calibrations.json; opis procedury: manuskrypt, sekcja 3.5 (część napisana przez autora rozprawy). Kodu projektu GHZ (katalog poza repozytorium klasyfikatora) nie czytano.
\agenttask{Pytanie do autora (recenzja rozdziału 8, ustalenie D14): według jakiej reguły kalibracja stanami GHZ dzieli zmierzoną stratę wierności na odstrojenie, tłumienie fazy i depolaryzację (trzy parametry dopasowane do jednej liczby; kod jest w projekcie GHZ poza repozytorium klasyfikatora) i co oznacza pole device_within_model = false we wszystkich sześciu wpisach pliku ibm_fez_ghz_calibrations.json? Po odpowiedzi zastąpić zdanie o braku opisu reguły jej krótkim opisem.}

\begin{table}[htb]
\centering
\caption[Modele klasyfikatora użyte w przebiegach na urządzeniu i w eksperymentach z korektą odczytu]{Modele klasyfikatora uruchomione na urządzeniu, w kolejności malejącej liczby bramek dwukubitowych CZ po transpilacji. $K$: liczba klas, $d$: liczba cech, $n$: liczba kubitów, $L$: liczba bloków; $\cohangle$ i $\qdep$: poziom ibm\_fez, z którym model uczono ($\qphase=0$, w modelu M10 $1{,}5\cdot10^{-4}$). ZP i S: udział w eksperymencie na zapisanych pomiarach i w eksperymencie symulacyjnym; (u) oznacza analizę uzupełniającą.}
\label{tab:cls:models}
\footnotesize
\setlength{\tabcolsep}{5pt}
\begin{tabular}{@{}llrrrrrrrcc@{}}
\toprule
Model & Identyfikator & $K$ & $d$ & $n$ & $L$ & CZ & $\cohangle$ [mrad] & $\qdep$ [$10^{-3}$] & ZP & S \\
\midrule
M1 & \texttt{c7\_f200\_q6\_b15} & 7 & 200 & 6 & 15 & 1814 & 1{,}28 & 2{,}98 & (u) & tak \\
M2 & \texttt{c7\_f200\_q3\_b13\_d2} & 7 & 200 & 3 & 13 & 233 & 1{,}28 & 2{,}98 & tak & tak \\
M3 & \texttt{c6\_f300\_q5\_b3} & 6 & 300 & 5 & 3 & 198 & 5{,}98 & 4{,}33 & tak & tak \\
M4 & \texttt{c8\_f200\_q7\_b1} & 8 & 200 & 7 & 1 & 138 & 5{,}66 & 3{,}63 & tak & -- \\
M5 & \texttt{c8\_f300\_q7\_b1} & 8 & 300 & 7 & 1 & 138 & 5{,}66 & 3{,}63 & tak & -- \\
M6 & \texttt{c7\_f200\_q6\_b1} & 7 & 200 & 6 & 1 & 113 & 1{,}28 & 2{,}98 & tak & tak \\
M7 & \texttt{c6\_f100\_q5\_b1} & 6 & 100 & 5 & 1 & 63 & 5{,}98 & 4{,}33 & tak & -- \\
M8 & \texttt{c6\_f300\_q5\_b1} & 6 & 300 & 5 & 1 & 63 & 5{,}98 & 4{,}33 & tak & -- \\
M9 & \texttt{c5\_f300\_q4\_b1} & 5 & 300 & 4 & 1 & 29 & 3{,}94 & 5{,}25 & tak & -- \\
M10 & \texttt{c4\_f400\_q3\_b1} & 4 & 400 & 3 & 1 & 15 & 3{,}70 & 3{,}81 & tak & tak \\
M11 & \texttt{c3\_f200\_q2\_b1} & 3 & 200 & 2 & 1 & 3 & 9{,}02 & 4{,}89 & tak & tak \\
\bottomrule
\end{tabular}
\end{table}
% ŹRÓDŁO tabeli: pliki experiment.json przebiegów testowych (config: n_classes, n_features, n_qubits, n_blocks, deltas_per_qubit, noise_strengths; transpile.two_qubit_gates) w katalogach wymienionych w eksperymenty/pdu-ibm-fez/klasyfikator/wyniki/metadane.json i punkty_odniesienia.json; udział w eksperymentach: pliki README.md obu eksperymentów.

\paragraph{Uczenie i dostrajanie na urządzeniu.}
Modele uczono ze współczynnikiem uczenia $0{,}001$, zmniejszanym o połowę po pięciu epokach bez poprawy straty testowej, na wsadach po $32$ próbki i z ziarnem $42$; punkty kontrolne pochodzą z epok od $16$ do $41$, w których dokładność na części uczącej wynosi już od $0{,}90$ do $0{,}99$.
% ŹRÓDŁO: /Users/krzysztofwerner/git/quantum-classifier/runs/<przebieg>/hyperparameters.json (learning_rate, batch_size, seed, scheduler, epochs) i metrics.json (best.epoch, train_accuracy) modeli M1--M11; hyperparameters.json przebiegów wsadowych (35--50 epok) i przeglądu szumu (35 epok); modele M1--M11: 50 albo 75 epok; eksperymenty/weryfikacja-liczb/rozdzial-8.csv, grupa H.
W każdej z trzech rund dostrajania mierzy się $400$ próbek uczących, wylosowanych warstwowo (jeśli zmierzona strata wzrosła względem ostatniego zaakceptowanego pomiaru, poprzednią aktualizację sieci cofa się, a współczynnik uczenia i liczbę epok zmniejsza o połowę; odrzucono $3$ z $20$ pomiarów), dopasowuje dla każdego kubitu metodą najmniejszych kwadratów odwzorowanie afiniczne wektora Blocha symulowanego w zmierzony (w pierwszej rundzie współczynnik determinacji $R^2$ od $0{,}91$ do $0{,}996$ w modelach M3--M11 i $0{,}78$ w modelu M2, więc odchylenie urządzenia od symulatora jest w dużej części systematyczne) i przez pięć epok uczy sieć na tych próbkach (Adam, współczynnik uczenia $10^{-4}$), licząc odczyt z obrazu wektora symulowanego w tym odwzorowaniu.
Po trzeciej rundzie ponownie mierzy się $100$ próbek testowych (tabela~\ref{tab:cls:hw-results}, omówienie w podrozdziale~\ref{subsec:cls:results:noise}).
% ŹRÓDŁO współczynnika determinacji: runs/<przebieg>/hwtrain_ibm_fez_*/experiment.json, history[0].device_model (r2_mean, r2_identity_mean; 400 próbek uczących; dla odwzorowania tożsamościowego R^2 od 0,99 w modelach M10 i M11 przez 0,38 w M4 do wartości ujemnych w M2 i M3).
% ŹRÓDŁO: torus_classifier/hardware_training.py (docstring modułu, funkcje acceptance, fit_device_model, train_epochs, prepare_setup; stałe DEFAULT_ROUNDS = 3, DEFAULT_EPOCHS_PER_ROUND = 5, DEFAULT_LR = 1e-4, DEFAULT_N_TRAIN = 400) i pola rounds, epochs_per_round, lr, train w plikach experiment.json dostrajania; odrzucone pomiary: history[].decision.accepted.

\begin{table}[htb]
\centering
\caption[Modele M1--M11 na symulatorze i na urządzeniu ibm\_fez przed dostrajaniem i po nim]{Modele M1--M11 na symulatorze i na urządzeniu ibm\_fez. CZ: liczba bramek dwukubitowych po transpilacji. Pełny zbiór: dokładność klasyfikatora (FBS, symulator) i sieci MLP o tych samych warstwach ukrytych na $2000$ próbkach testowych. S: symulator z szumem uczenia (wartości dokładne), U: urządzenie; zg.: zgodność predykcji obu źródeł; S/U: liczba próbek poprawnych tylko na symulatorze i tylko na urządzeniu; $+/-$: liczba próbek, które po dostrajaniu stały się poprawne i błędne; $p$: $p$-wartość dokładnego testu McNemara.}
\label{tab:cls:hw-results}
\footnotesize
\setlength{\tabcolsep}{3.6pt}
\begin{tabular}{@{}lrcccccccccccc@{}}
\toprule
 & & \multicolumn{2}{c}{Pełny zbiór} & \multicolumn{5}{c}{Przebieg testowy, $100$ próbek} & \multicolumn{2}{c}{Średni promień} & \multicolumn{3}{c}{Po dostrajaniu} \\
\cmidrule(lr){3-4}\cmidrule(lr){5-9}\cmidrule(lr){10-11}\cmidrule(l){12-14}
Model & CZ & FBS & MLP & S & U & zg. & S/U & $p$ & S & U & U & $+/-$ & $p$ \\
\midrule
M1 & 1814 & $0{,}378$ & -- & $0{,}34$ & $0{,}13$ & $0{,}12$ & 29/8 & $<0{,}001$ & $0{,}69$ & $0{,}09$ & -- & -- & -- \\
M2 & 233 & $0{,}439$ & $0{,}520$ & $0{,}45$ & $0{,}35$ & $0{,}71$ & 14/4 & $0{,}03$ & $0{,}84$ & $0{,}37$ & $0{,}41$ & 11/5 & $0{,}21$ \\
M3 & 198 & $0{,}455$ & $0{,}509$ & $0{,}46$ & $0{,}43$ & $0{,}91$ & 4/1 & $0{,}38$ & $0{,}90$ & $0{,}50$ & $0{,}46$ & 3/0 & $0{,}25$ \\
M4 & 138 & $0{,}425$ & $0{,}477$ & $0{,}41$ & $0{,}42$ & $0{,}97$ & 0/1 & $1$ & $0{,}92$ & $0{,}67$ & $0{,}46$ & 6/2 & $0{,}29$ \\
M5 & 138 & $0{,}327$ & $0{,}400$ & $0{,}43$ & $0{,}41$ & $0{,}96$ & 2/0 & $0{,}50$ & $0{,}88$ & $0{,}64$ & $0{,}39$ & 3/5 & $0{,}73$ \\
M6 & 113 & $0{,}455$ & $0{,}520$ & $0{,}46$ & $0{,}47$ & $0{,}89$ & 1/2 & $1$ & $0{,}91$ & $0{,}61$ & $0{,}48$ & 5/4 & $1$ \\
M7 & 63 & $0{,}606$ & $0{,}660$ & $0{,}53$ & $0{,}53$ & $0{,}97$ & 1/1 & $1$ & $0{,}92$ & $0{,}71$ & $0{,}51$ & 5/7 & $0{,}77$ \\
M8 & 63 & $0{,}503$ & $0{,}524$ & $0{,}51$ & $0{,}50$ & $0{,}99$ & 1/0 & $1$ & $0{,}95$ & $0{,}77$ & $0{,}49$ & 3/4 & $1$ \\
M9 & 29 & $0{,}568$ & $0{,}605$ & $0{,}62$ & $0{,}61$ & $0{,}98$ & 1/0 & $1$ & $0{,}95$ & $0{,}78$ & $0{,}58$ & 3/6 & $0{,}51$ \\
M10 & 15 & $0{,}624$ & $0{,}648$ & $0{,}58$ & $0{,}58$ & $1{,}00$ & 0/0 & $1$ & $0{,}88$ & $0{,}83$ & $0{,}60$ & 6/4 & $0{,}75$ \\
M11 & 3 & $0{,}807$ & $0{,}818$ & $0{,}82$ & $0{,}82$ & $1{,}00$ & 0/0 & $1$ & $0{,}99$ & $0{,}94$ & $0{,}81$ & 0/1 & $1$ \\
\bottomrule
\end{tabular}
\end{table}
% ŹRÓDŁO tabeli: przebiegi testowe runs/<przebieg>/hardware_ibm_fez_*/ modeli M1--M11: experiment.json (transpile.two_qubit_gates, checkpoint.test_accuracy, mlp.best_test_accuracy), summary.json (sources.sim_noisy.r_mean, sources.hardware.r_mean), results.csv (kolumny sim_noisy_correct, hardware_correct, sim_noisy_pred, hardware_pred); dostrajanie: runs/<przebieg>/hwtrain_ibm_fez_*/eval/final_hardware/results.csv (hardware_correct). Katalogi przebiegów: eksperymenty/weryfikacja-liczb/rozdzial-8-skrypty/common.py. Wiersze wygenerował skrypt rysunki.py; testy McNemara policzono w rozprawie.

\paragraph{Korekta odczytu i eksperyment porównawczy.}
Tabela~\ref{tab:cls:pdu-recovery} uzupełnia podrozdział~\ref{subsec:cls:results:pdu}, a tabela~\ref{tab:cls:app:family} podaje pięć porównań zaplanowanych w każdej z dwóch konfiguracji eksperymentu porównawczego (podrozdział~\ref{subsec:cls:results:training}).

\begin{table}[htb]
\centering
\caption[Korekta odczytu na progu spadku dokładności w eksperymencie symulacyjnym]{Korekta odczytu na progu spadku dokładności w eksperymencie symulacyjnym (modele wieloblokowe, $\nshots=1024$). Próg: siła składowej, przy której dokładność bez korekty spada o $0{,}05$; poziom: pierwszy poziom siatki, na którym spadek osiąga $0{,}05$, i średni promień wektora Blocha na tym poziomie. Dokładność bez szumu obwodu wynosi $0{,}377$ (M1), $0{,}437$ (M2) i $0{,}453$ (M3); w nawiasach odzysk w procentach.}
\label{tab:cls:pdu-recovery}
\footnotesize
\setlength{\tabcolsep}{4pt}
\begin{tabular}{@{}llrrrcccc@{}}
\toprule
 & & & & & & \multicolumn{3}{c}{Dokładność po korekcie (odzysk)} \\
\cmidrule(l){7-9}
Model & Składowa & Próg & Poziom & Promień & Bez korekty & FZ & AO & AR \\
\midrule
M1 & odstrojenie & $0{,}107$ & $0{,}2$ & $0{,}30$ & $0{,}142$ & $0{,}218$ ($32$) & $0{,}211$ ($29$) & $0{,}319$ ($75$) \\
 & tłumienie fazy & $0{,}146$ & $0{,}2$ & $0{,}25$ & $0{,}285$ & $0{,}352$ ($73$) & $0{,}312$ ($29$) & $0{,}363$ ($85$) \\
 & depolaryzacja & $0{,}040$ & $0{,}05$ & $0{,}09$ & $0{,}290$ & $0{,}318$ ($32$) & $0{,}241$ ($-57$) & $0{,}304$ ($16$) \\
\addlinespace
M2 & odstrojenie & $0{,}067$ & $0{,}07$ & $0{,}85$ & $0{,}378$ & $0{,}384$ ($11$) & $0{,}421$ ($74$) & $0{,}427$ ($83$) \\
 & tłumienie fazy & $0{,}186$ & $0{,}2$ & $0{,}41$ & $0{,}383$ & $0{,}383$ ($-1$) & $0{,}416$ ($61$) & $0{,}412$ ($53$) \\
 & depolaryzacja & $0{,}071$ & $0{,}1$ & $0{,}07$ & $0{,}279$ & $0{,}301$ ($14$) & $0{,}220$ ($-37$) & $0{,}262$ ($-11$) \\
\addlinespace
M3 & odstrojenie & $0{,}237$ & $0{,}3$ & $0{,}90$ & $0{,}329$ & $0{,}366$ ($30$) & $0{,}454$ ($100$) & $0{,}453$ ($100$) \\
 & tłumienie fazy & \multicolumn{7}{l}{próg nieosiągnięty: do siły $0{,}5$ zmiana dokładności poniżej $0{,}002$} \\
 & depolaryzacja & $0{,}309$ & $0{,}5$ & $0{,}01$ & $0{,}192$ & $0{,}155$ ($-14$) & $0{,}181$ ($-4$) & $0{,}154$ ($-14$) \\
\bottomrule
\end{tabular}
\end{table}
% ŹRÓDŁO tabeli: eksperymenty/krzywe-degradacji/klasyfikator/wyniki/progi.csv (2000 próbek testowych, średnie z dziesięciu powtórzeń; siła progowa z interpolacji liniowej w logarytmie siły), wiersze o strzaly = 1024 modeli c7_f200_q6_b15 (M1), c7_f200_q3_b13_d2 (M2) i c6_f300_q5_b3 (M3): sila_progowa, poziom_siatki, promien_analityczny, dokl_zero_szumu, dokl_bez_korekty, dokl_a2 i odzysk_a2, dokl_b i odzysk_b, dokl_b2 i odzysk_b2; te same wartości z większą liczbą cyfr w pliku krzywe.csv (dokl, odzysk_dokl); wiersz M3, tłumienie fazy: progi.csv, spadek_przy_sile_maks = -0,0016. Zaokrąglenia połówek w górę (na przykład 0,31885 to 0,319). Pominięte dla budżetu stron wiersze przebiegu mieszanego (siła = krotność poziomu ibm_fez; próg, poziom, promień, bez korekty, FZ, AO, AR z odzyskiem w procentach): M1 12,7; 20; 0,06; 0,234; 0,260 (18); 0,201 (-23); 0,238 (3); M2 22,6; 50; 0,02; 0,169; 0,160 (-3); 0,165 (-1); 0,155 (-5); M3 28,6; 50; 0,19; 0,309; 0,378 (47); 0,446 (95); 0,452 (99).

\begin{table}[htb]
\centering
\caption[Porównania rodziny głównej eksperymentu porównawczego]{Porównania rodziny głównej eksperymentu porównawczego ($\nshots=1024$, pięć ziaren): średnia różnica dokładności sparowana po ziarnach, $95$-procentowy przedział ufności z rozkładu t~Studenta i $p$-wartość testu t po poprawce Holma w rodzinie pięciu porównań danej konfiguracji ($p_{\mathrm{H}}$). Konfiguracja płytka: szum testu na poziomie $0{,}10$; głęboka (model M2): na poziomie $0{,}05$. FZ i AR: model uczony bez szumu z korektą odczytu.}
\label{tab:cls:app:family}
\footnotesize
\setlength{\tabcolsep}{5pt}
\begin{tabular}{@{}lcccc@{}}
\toprule
 & \multicolumn{2}{c}{Konfiguracja płytka} & \multicolumn{2}{c}{Konfiguracja głęboka} \\
\cmidrule(lr){2-3}\cmidrule(l){4-5}
Porównanie & różnica & $p_{\mathrm{H}}$ & różnica & $p_{\mathrm{H}}$ \\
\midrule
trening w szumie $-$ bez działań & $+0{,}046$ $[-0{,}070;\,+0{,}162]$ & $1$ & $+0{,}059$ $[-0{,}025;\,+0{,}143]$ & $0{,}45$ \\
FZ $-$ bez działań & $+0{,}031$ $[-0{,}045;\,+0{,}107]$ & $1$ & $-0{,}003$ $[-0{,}021;\,+0{,}014]$ & $0{,}62$ \\
AR $-$ bez działań & $+0{,}049$ $[-0{,}059;\,+0{,}156]$ & $1$ & $+0{,}026$ $[-0{,}025;\,+0{,}076]$ & $0{,}46$ \\
trening w szumie $-$ FZ & $+0{,}015$ $[-0{,}036;\,+0{,}066]$ & $1$ & $+0{,}062$ $[-0{,}016;\,+0{,}140]$ & $0{,}45$ \\
trening w szumie $-$ AR & $-0{,}003$ $[-0{,}047;\,+0{,}041]$ & $1$ & $+0{,}033$ $[-0{,}008;\,+0{,}075]$ & $0{,}45$ \\
\bottomrule
\end{tabular}
\end{table}
% ŹRÓDŁO tabeli: eksperymenty/trening-w-szumie/wyniki/porownania.csv (rodzina glowna, tryb strzaly, kalibracja uczaca, E3) i eksperymenty/trening-w-szumie/gleboka/wyniki/porownania.csv (rodzina glowna, E2): kolumny roznica, ci95_lo, ci95_hi, p_t_holm_w_rodzinie_glownej; wiersze kolejno: wariant 3 - wariant 1, wariant 2 a2 - wariant 1, wariant 2 b2 - wariant 1, wariant 3 - wariant 2 a2, wariant 3 - wariant 2 b2. Wartości niezaokrąglone, płytka: +0,0459 [-0,0704; +0,1622]; +0,0310 [-0,0450; +0,1070]; +0,0488 [-0,0586; +0,1563]; +0,0149 [-0,0357; +0,0655]; -0,0029 [-0,0468; +0,0409]; p Holma 1 w pięciu wierszach (p bez poprawki 0,33; 0,32; 0,28; 0,46; 0,86). Głęboka: +0,0588 [-0,0255; +0,1432]; -0,0035 [-0,0214; +0,0144]; +0,0255 [-0,0246; +0,0757]; +0,0623 [-0,0156; +0,1402]; +0,0333 [-0,0083; +0,0749]; p Holma 0,452; 0,620; 0,461; 0,452; 0,452 (bez poprawki 0,12; 0,62; 0,23; 0,09; 0,09).
